\documentclass{article}
\usepackage{xeCJK}
\begin{document}
% 验和测试数据注释行，包含中文标点，。and a trailing comment with 50% literal
通过实验验证，准确率达到 50\% 的提升 and R\&D notes here 和后续分析。
我们测试了 50% off 的折扣 and the Smith & Jones method 在数据集上的表现。
第一句结束。  第二句开始 with a double space after the period 和中文。
如图所示 \cite{zhang2024} 的方法 and another result \cite{li2025} 很有效。
参数设置为 a ; b 以及 c : d 还有 e~~f 的非断行空格 在中文语境中。
使用 "straight quotes" 和 \'apostrophes\' 在 ASCII context 验和测试。
范围是 5 -- 7 以及 spaced 5 – 3 的减号用法 验和数据 2020–2025。
内联公式 $x = 验 + 和$ 以及 $a \times b \div c => d$ 的数学表达式。
\begin{equation} 验和 = \alpha , \beta . \end{equation}
全角标点测试，。：；！？ in ASCII context 验和 with fullwidth punctuation。
省略号用法 ... 以及 trailing spaces below 验和   
% 第二条注释 with CJK 验和 and a mid-comment 75% value 不应被转义
末尾段落 验和测试 final paragraph with mixed 中英文 content and 100% coverage.
\end{document}
